\subsubsection{Fortran Utils}
\begin{enumerate}
    \item \textbf{Binary Search Trees for 1D Range Queries}
          
          Binary Search Trees provide efficient indexing and range query capabilities
          for one-dimensional data using flat array implementations for optimal memory
          performance.
          
          \textbf{Fortran BST Implementation}
          
          \begin{enumerate}
            \item \textbf{Build BST Index}
                  
                  Constructs a Binary Search Tree index by sorting the input values and storing
                  permutation indices. Uses an efficient stack-based sorting algorithm.
                  
                  \textbf{Input:}
                  \begin{itemize}
                    \item \texttt{values real(real64)}: 1D array of values to index
                          
                    \item \texttt{num\_values integer(int32)}: Number of elements
                  \end{itemize}
                  \textbf{Output:}
                  \begin{itemize}
                    \item \texttt{sorted\_indices integer(int32)}: Permutation indices after
                          sorting (1-based)
                          
                    \item \texttt{left\_stack, right\_stack integer(int32)}: Sorting workspace
                          
                    \item \texttt{ierr integer(int32)}: Error code
                  \end{itemize}
                  \begin{lstlisting}[language=Fortran]
call build_bst_index(values, num_values, sorted_indices, left_stack, &
                        right_stack, ierr)
                  \end{lstlisting}
                  
            \item \textbf{Get Sorted Value}
                  
                  Retrieves the value at a specific position in the sorted order defined
                  by the BST index.
                  
                  \textbf{Input:}
                  \begin{itemize}
                    \item \texttt{values real(real64)}: Original values array
                          
                    \item \texttt{sorted\_indices integer(int32)}: BST index array (1-based)
                          
                    \item \texttt{position integer(int32)}: Position in sorted order (1-based)
                  \end{itemize}
                  \textbf{Output:}
                  \begin{itemize}
                    \item \texttt{sorted\_value real(real64)}: Value at specified position
                          
                    \item \texttt{ierr integer(int32)}: Error code
                  \end{itemize}
                  \begin{lstlisting}[language=Fortran]
sorted_value = get_sorted_value(values, sorted_indices, &
                                position, ierr)
                  \end{lstlisting}
                  
            \item \textbf{BST Range Query}
                  
                  Efficiently finds all values within a specified range using the precomputed
                  BST index.
                  
                  \textbf{Input:}
                  \begin{itemize}
                    \item \texttt{values real(real64)}: Original values array
                          
                    \item \texttt{sorted\_indices integer(int32)}: BST index array (1-based)
                          
                    \item \texttt{num\_values integer(int32)}: Number of elements
                          
                    \item \texttt{lower\_bound, upper\_bound real(real64)}: Range boundaries
                  \end{itemize}
                  \textbf{Output:}
                  \begin{itemize}
                    \item \texttt{output\_indices integer(int32)}: Indices of matching values
                          (1-based)
                          
                    \item \texttt{num\_matches integer(int32)}: Number of matches found
                          
                    \item \texttt{ierr integer(int32)}: Error code
                  \end{itemize}
                  \begin{lstlisting}[language=Fortran]
call bst_range_query(values, sorted_indices, num_values, &
                    lower_bound, upper_bound, output_indices, &
                    num_matches, ierr)
                  \end{lstlisting}
          \end{enumerate}
          
    \item \textbf{k-d Trees for Multi-Dimensional Indexing} \label{sec:kd_trees}
          
          k-d trees provide efficient space partitioning for multi-dimensional data with
          support for both Euclidean and spherical geometries.
          
          \textbf{Fortran k-d Tree Implementation}
          
          \begin{enumerate}
            \item \textbf{Build k-d Tree Index}
                  
                  Constructs a k-d tree index using stack-based non-recursive partitioning
                  along specified dimensions.
                  
                  \textbf{Input:}
                  \begin{itemize}
                    \item \texttt{points real(real64)}: Data points matrix (dimensions ×
                          points)
                          
                    \item \texttt{num\_dimensions integer(int32)}: Number of dimensions
                          
                    \item \texttt{num\_points integer(int32)}: Number of points
                          
                    \item \texttt{dimension\_order integer(int32)}: Dimension ordering for
                          splitting (1-based)
                  \end{itemize}
                  \textbf{Output:}
                  \begin{itemize}
                    \item \texttt{kd\_indices integer(int32)}: Output index array (1-based)
                          
                    \item \texttt{workspace integer(int32)}: Workspace array
                          
                    \item \texttt{value\_buffer real(real64)}: Value buffer for sorting
                          
                    \item \texttt{permutation integer(int32)}: Permutation array
                          
                    \item \texttt{left\_stack, right\_stack integer(int32)}: Sorting stacks
                          
                    \item \texttt{recursion\_stack integer(int32)}: Recursion stack
                          
                    \item \texttt{ierr integer(int32)}: Error code
                  \end{itemize}
                  \begin{lstlisting}[language=Fortran]
call build_kd_index(points, num_dimensions, num_points, kd_indices, &
                    dimension_order, workspace, value_buffer, &
                    permutation, left_stack, right_stack, &
                    recursion_stack, ierr)
                  \end{lstlisting}
                  
            \item \textbf{Build Spherical k-d Tree}
                  
                  Specialized k-d tree construction optimized for unit vectors on hyperspheres.
                  
                  \textbf{Input:}
                  \begin{itemize}
                    \item \texttt{vectors real(real64)}: Unit vectors matrix (dimensions
                          × vectors)
                          
                    \item \texttt{num\_dimensions integer(int32)}: Number of dimensions
                          
                    \item \texttt{num\_vectors integer(int32)}: Number of vectors
                          
                    \item \texttt{dimension\_order integer(int32)}: Dimension ordering for
                          splitting (1-based)
                  \end{itemize}
                  \textbf{Output:}
                  \begin{itemize}
                    \item \texttt{sphere\_indices integer(int32)}: Output index array (1-based)
                          
                    \item Various workspace arrays (same as build\_kd\_index)
                          
                    \item \texttt{ierr integer(int32)}: Error code
                  \end{itemize}
                  \begin{lstlisting}[language=Fortran]
call build_spherical_kd(vectors, num_dimensions, num_vectors, &
                        sphere_indices, dimension_order, workspace, &
                        value_buffer, permutation, left_stack, & 
                        right_stack, recursion_stack, ierr)
                  \end{lstlisting}
                  
            \item \textbf{Get Point from k-d Index}
                  
                  Retrieves point coordinates from specified position in k-d tree ordering.
                  
                  \textbf{Input:}
                  \begin{itemize}
                    \item \texttt{points real(real64)}: Original points matrix
                          
                    \item \texttt{kd\_indices integer(int32)}: k-d tree index array (1-based)
                          
                    \item \texttt{position integer(int32)}: Position in index (1-based)
                  \end{itemize}
                  \textbf{Output:}
                  \begin{itemize}
                    \item \texttt{point\_values real(real64)}: Retrieved point values
                          
                    \item \texttt{ierr integer(int32)}: Error code
                  \end{itemize}
                  \begin{lstlisting}[language=Fortran]
call get_kd_point(points, kd_indices, position, point_values, ierr)
                  \end{lstlisting}
          \end{enumerate}
          \item \textbf{Hashing}
          
          The TensorOmics package provides high-performance hashing utilities using the
          XXH3 algorithm for both hash maps and hash sets. These utilities are
          implemented in Fortran with separate chaining collision resolution and
          automatic resizing.
          
          \begin{enumerate}
            \item \textbf{Hashmap}
                  
                  The hashmap implementation provides key-value storage with string keys
                  and integer values, using the XXH3 hashing algorithm for optimal
                  performance.
                  
                  \begin{itemize}
                    \item \textbf{Type Definition:}
                          \begin{lstlisting}[language=Fortran]
type(hashmap_type) :: map
                          \end{lstlisting}
                          
                    \item \textbf{Create Hashmap:} Initialize a new hashmap with optional initial size. The Hashmap size will be set to the next greater
                          power of two.
                          \begin{lstlisting}[language=Fortran]
call hashmap_create(map, initial_size)
                          \end{lstlisting}
                          \begin{itemize}
                            \item \texttt{map}: Hashmap object to create
                                  
                            \item \textbf{(Optional)} \texttt{initial\_size}: Initial size
                                  of the hashmap
                          \end{itemize}
                          
                    \item \textbf{Insert Key-Value Pair:} Insert or update a key-value pair in the hashmap. 
                          \begin{lstlisting}[language=Fortran]
call hashmap_put(map, key, value)
                    \end{lstlisting}
                    \begin{itemize}
                        \item \texttt{map}: Hashmap to insert into
                              
                        \item \texttt{key}: String key to store
                              
                        \item \texttt{value}: Integer value to store
                    \end{itemize}
                    
                    \item \textbf{Lookup Key:} Retrieve the value associated with a key.
                          \begin{lstlisting}[language=Fortran]
value = hashmap_get(map, key)
                          \end{lstlisting}
                          \begin{itemize}
                            \item \texttt{map}: Hashmap to search
                                  
                            \item \texttt{key}: String key to look for
                                  
                            \item \textbf{Returns:} Integer value, or -1 if key not found
                          \end{itemize}
                          
                    \item \textbf{Destroy Hashmap:} Clean up and deallocate the hashmap.
                          It is important to always call this function when the hashmap is no
                          longer in use to avoid memory leaks.
                          \begin{lstlisting}[language=Fortran]
call hashmap_destroy(map)
                          \end{lstlisting}
                          \begin{itemize}
                            \item \texttt{map}: Hashmap to destroy
                          \end{itemize}
                  \end{itemize}
                  
                  \textbf{Hashmap Usage Example:}
                  \begin{lstlisting}[language=Fortran]
use xxh3_hashmap_module

type(hashmap_type) :: gene_map
integer :: gene_index

! Create hashmap for gene IDs
call hashmap_create(gene_map, initial_size=1000)

! Store gene indices
call hashmap_put(gene_map, "Gene_001", 1)
call hashmap_put(gene_map, "Gene_002", 2)
call hashmap_put(gene_map, "Gene_003", 3)

! Lookup gene index
gene_index = hashmap_get(gene_map, "Gene_002")
! gene_index = 2

! Clean up
call hashmap_destroy(gene_map)
                  \end{lstlisting}
                  
            \item \textbf{Hashset}
                  
                  The hashset implementation provides unique string storage with fast membership
                  testing, using the same XXH3 hashing algorithm as the hashmap.
                  
                  \begin{itemize}
                    \item \textbf{Type Definition:}
                          \begin{lstlisting}[language=Fortran]
type(hashset_type) :: set
                          \end{lstlisting}
                          
                    \item \textbf{Create Hashset:} Initialize a new hashset with
                          optional initial size. The hashset size will be increased to the next
                          greatest power of two.
                          \begin{lstlisting}[language=Fortran]
call hashset_create(set, initial_size)
                          \end{lstlisting}
                          \begin{itemize}
                            \item \texttt{set}: Hashset object to create
                                  
                            \item \textbf{(Optional)} \texttt{initial\_size}: Initial size
                                  of the hashset
                          \end{itemize}
                          
                    \item \textbf{Insert Key:} Insert a unique key into the hashset. 
                    \begin{lstlisting}[language=Fortran]
call hashset_put(set, key, ierr)
                    \end{lstlisting}
                    \begin{itemize}
                        \item \texttt{set}: Hashset to insert into
                              
                        \item \texttt{key}: String key to store
                              
                        \item \texttt{ierr}: Error code (0 if successful, error if duplicate)
                    \end{itemize}
                    
                    \item \textbf{Check Membership:} Test if a key exists in the hashset.
                          \begin{lstlisting}[language=Fortran]
exists = is_in_hashset(set, key)
                          \end{lstlisting}
                          \begin{itemize}
                            \item \texttt{set}: Hashset to search
                                  
                            \item \texttt{key}: String key to check
                                  
                            \item \textbf{Returns:} Logical indicating membership
                          \end{itemize}
                          
                    \item \textbf{Destroy Hashset:} Clean up and deallocate the hashset.
                          It is important that this function is always called when the hashset
                          is no longer in use to avoid memory leaks.
                          \begin{lstlisting}[language=Fortran]
call hashset_destroy(set)
                          \end{lstlisting}
                          \begin{itemize}
                            \item \texttt{set}: Hashset to destroy
                          \end{itemize}
                  \end{itemize}
                  
                  \textbf{Hashset Usage Example:}
                  \begin{lstlisting}[language=Fortran]
use xxh3_hashmap_module

type(hashset_type) :: unique_genes
integer :: ierr
logical :: exists

! Create hashset for unique gene validation
call hashset_create(unique_genes, initial_size=5000)

! Add genes to set
call hashset_put(unique_genes, "Gene_001", ierr)
call hashset_put(unique_genes, "Gene_002", ierr)
call hashset_put(unique_genes, "Gene_003", ierr)

! Check for duplicates (will return error)
call hashset_put(unique_genes, "Gene_001", ierr)
! ierr != 0 indicates duplicate key

! Test membership
exists = is_in_hashset(unique_genes, "Gene_002")
! exists = .true.

! Clean up
call hashset_destroy(unique_genes)
                  \end{lstlisting}
          \end{enumerate}
          
          \textbf{Technical Features:}
          \begin{itemize}
            \item Uses XXH3 non-cryptographic hash algorithm for optimal performance
                  
            \item Automatic resizing when load factor exceeds 0.75
                  
            \item Separate chaining collision resolution
                  
            \item Table sizes are powers of two for efficient modulo operations
                  
            \item Keys are automatically trimmed and normalized
                  
            \item Memory efficient with proper cleanup procedures
                  
            \item Thread-safe for read operations (concurrent writes not supported)
          \end{itemize}
          
          \textbf{Performance Characteristics:}
          \begin{itemize}
            \item Average case O(1) for insertions and lookups
                  
            \item Memory usage proportional to number of elements + overhead
                  
            \item Resizing operations are O(n) but amortized over many insertions
          \end{itemize}
          
    \item \textbf{Sorting Algorithms}
          
          The TensorOmics package provides high-performance quicksort implementations
          for various data types using iterative algorithms with manual stack management.
          These sorting routines operate indirectly through permutation vectors to
          preserve the original data order.
          
          \begin{enumerate}
            \item \textbf{Generic Sorting Interface}
                  
                  A unified interface for sorting different data types through the \texttt{sort\_array}
                  generic procedure.
                  
                  \begin{itemize}
                    \item \textbf{Generic Procedure:}
                          \begin{lstlisting}[language=Fortran]
call sort_array(array, perm, stack_left, stack_right)
                          \end{lstlisting}
                          \begin{itemize}
                            \item \texttt{array}: Input array to sort (real, integer, or
                                  character)
                                  
                            \item \texttt{perm}: Permutation vector that will be reordered
                                  
                            \item \texttt{stack\_left}: Work array for left indices (size =
                                  array size)
                                  
                            \item \texttt{stack\_right}: Work array for right indices (size
                                  = array size)
                          \end{itemize}
                          
                          \textbf{Supported Types:}
                          \begin{lstlisting}[language=Fortran]
! Real arrays
real(real64) :: data_real(:)
! Integer arrays  
integer(int32) :: data_int(:)
! Character arrays
character(len=*) :: data_char(:)
                          \end{lstlisting}
                  \end{itemize}
                  
            \item \textbf{Type-Specific Sorting Procedures}
                  
                  Direct access to type-specific sorting routines for explicit control.
                  
                  \begin{itemize}
                    \item \textbf{Sort Real Arrays:}
                          \begin{lstlisting}[language=Fortran]
call sort_real(array, perm, stack_left, stack_right)
                          \end{lstlisting}
                          \begin{itemize}
                            \item \texttt{array}: Real(real64) input array to sort
                                  
                            \item \texttt{perm}: Permutation vector that will be sorted
                                  
                            \item \texttt{stack\_left}: Manual stack of left indices
                                  
                            \item \texttt{stack\_right}: Manual stack of right indices
                          \end{itemize}
                          
                    \item \textbf{Sort Integer Arrays:}
                          \begin{lstlisting}[language=Fortran]
call sort_integer(array, perm, stack_left, stack_right)
                          \end{lstlisting}
                          \begin{itemize}
                            \item \texttt{array}: Integer(int32) input array to sort
                                  
                            \item \texttt{perm}: Permutation vector that will be sorted
                                  
                            \item \texttt{stack\_left}: Manual stack of left indices
                                  
                            \item \texttt{stack\_right}: Manual stack of right indices
                          \end{itemize}
                          
                    \item \textbf{Sort Character Arrays:}
                          \begin{lstlisting}[language=Fortran]
call sort_character(array, perm, stack_left, stack_right)
                          \end{lstlisting}
                          \begin{itemize}
                            \item \texttt{array}: Character input array to sort (lexicographic
                                  order)
                                  
                            \item \texttt{perm}: Permutation vector that will be sorted
                                  
                            \item \texttt{stack\_left}: Manual stack of left indices
                                  
                            \item \texttt{stack\_right}: Manual stack of right indices
                          \end{itemize}
                  \end{itemize}
                  
                  \textbf{Sorting Usage Example:} 
                  \begin{lstlisting}[language=Fortran]
use f42_utils
use iso_fortran_env, only: real64, int32

real(real64) :: expression_data(1000)
integer(int32) :: perm(1000), stack_left(1000), stack_right(1000)
integer :: i

! Initialize permutation
do i = 1, 1000
perm(i) = i
end do

! Sort expression data indirectly
call sort_real(expression_data, perm, stack_left, stack_right)

! Access sorted data through permutation
! expression_data(perm(1)) is the smallest value
! expression_data(perm(1000)) is the largest value
            \end{lstlisting}
            
            \item \textbf{Technical Implementation Details}
                  
                  \begin{itemize}
                    \item \textbf{Algorithm:} Iterative quicksort with manual stack
                          management
                          
                    \item \textbf{Sorting Method:} Indirect sorting via permutation
                          vectors
                          
                    \item \textbf{Pivot Selection:} Median-of-three (center element)
                          
                    \item \textbf{Memory:} Requires two work arrays of same size as
                          input
                          
                    \item \textbf{Performance:} Average case O(n log n), worst case O(n²)
                          
                    \item \textbf{Stability:} Not stable (order of equal elements not preserved)
                  \end{itemize}
                  
            \item \textbf{Work Array Requirements}
                  
                  For sorting an array of size \texttt{n}, you must provide:
                  \begin{itemize}
                    \item \texttt{perm(n)}: Permutation vector (initialize with 1, 2, ...,
                          n)
                          
                    \item \texttt{stack\_left(n)}: Work array for stack operations
                          
                    \item \texttt{stack\_right(n)}: Work array for stack operations
                  \end{itemize}
                  
                  \textbf{Initialization Pattern:}
                  \begin{lstlisting}[language=Fortran]
do i = 1, n
    perm(i) = i          ! Identity permutation
end do
                  \end{lstlisting}
                  
            \item \textbf{Advanced Usage: Direct Algorithm Access}
                  
                  For maximum performance in tight loops, you can call the internal
                  quicksort procedures directly:
                  
                  \begin{lstlisting}[language=Fortran]
! Internal quicksort for real arrays
call quicksort_real(array, perm, n, stack_left, stack_right)

! Internal quicksort for integer arrays  
call quicksort_int(array, perm, n, stack_left, stack_right)

! Internal quicksort for character arrays
call quicksort_char(array, perm, n, stack_left, stack_right)
                  \end{lstlisting}
                  
                  \begin{itemize}
                    \item These procedures require explicit array size parameter \texttt{n}
                          
                    \item Work arrays must be properly initialized
                          
                    \item Provides slightly better performance by skipping interface overhead
                  \end{itemize}
          \end{enumerate}
          
          \textbf{Performance Considerations:}
          \begin{itemize}
            \item Manual stack management avoids recursion limits for large arrays
                  
            \item Indirect sorting preserves original data and allows multiple sort orders
                  
            \item Character sorting uses Fortran's intrinsic lexicographic comparison
                  
            \item Memory usage is $4n$ for all data including input, work array, and
                  both stacks.
                  
            \item In case of memory considerations, heapsort might be used.
          \end{itemize}
          
          \textbf{Error Handling:}
          \begin{itemize}
            \item All sorting procedures are \texttt{pure} (no side effects)
                  
            \item No explicit error codes - assumes valid input dimensions
                  
            \item User must ensure work arrays are properly sized
                  
            \item Invalid permutations may result in incorrect sorting
          \end{itemize}
\end{enumerate}